\documentclass[11pt]{article}

\usepackage[margin=1in]{geometry}
\usepackage{amsmath, amssymb, amsthm}
\usepackage{enumitem}

\pagestyle{empty}

\begin{document}

\begin{center}
  {\Large\bfseries Weekly Assignment: Division into Cases,\\[2pt]
   Quotient-Remainder, \& Floor/Ceiling\\[2pt]
   (Epp~4.5 \& 4.6)}
\end{center}

\bigskip

\noindent
Complete each of the following problems.  Show all work and justify your answers.

\bigskip

%% ─────────────────────────────────────────────
\noindent\textbf{Problem 1} (Epp 4.5 \#17). \;
Prove directly from the definitions that for every integer~$n$,
$\;n^2 - n + 3\;$ is odd.  Use division into two cases: $n$~is
even and $n$~is odd.

\bigskip

%% ─────────────────────────────────────────────
\noindent\textbf{Problem 2} (Epp 4.5 \#18\,a). \;
Prove that the product of any two consecutive integers is even.

\bigskip

%% ─────────────────────────────────────────────
\noindent\textbf{Problem 3} (Epp 4.5 \#40). \;
Prove that if $n$ is any odd integer, then
$\;n^4 \bmod 16 = 1$.

\medskip\noindent
\emph{Hint}: Recall that an odd number can be written as
$n = 2k + 1$ for some integer~$k$.  Expand $(2k+1)^4$ and
simplify modulo~16.

\bigskip

%% ─────────────────────────────────────────────
\noindent\textbf{Problem 4} (Epp 4.6 \#11). \;
State a necessary and sufficient condition for the floor of a real
number to equal that number.

\bigskip

%% ─────────────────────────────────────────────
\noindent\textbf{Problem 5} (Proof: Divisibility by 3). \;
Prove the following well-known divisibility rule from grade
school:

\begin{quote}
  \emph{If the sum of the digits of a (positive) integer is
  divisible by~$3$, then that integer is divisible by~$3$.}
\end{quote}

\noindent
\textbf{Hints}:
\begin{enumerate}[itemsep=4pt]
  \item A number written in decimal is a sum of its digits
        multiplied by powers of~10.  That is, a number with
        digits $d_k\, d_{k-1}\, \cdots\, d_1\, d_0$ equals
        \[
          d_k \cdot 10^k + d_{k-1} \cdot 10^{k-1} + \cdots
          + d_1 \cdot 10 + d_0.
        \]
  \item Recall that $n$ is divisible by~3 when
        $n \bmod 3 = 0$.
  \item The following properties of modular arithmetic may
        be useful:
        \begin{align*}
          (n + m) \bmod q &= \bigl((n \bmod q) +
                             (m \bmod q)\bigr) \bmod q, \\
          (n \cdot m) \bmod q &= \bigl((n \bmod q) \cdot
                             (m \bmod q)\bigr) \bmod q.
        \end{align*}
  \item Observe that $10 \bmod 3 = 1$, so
        $10^j \bmod 3 = 1$ for every $j \ge 0$.
\end{enumerate}

\end{document}
